\chapter{周期势与电子传播}\label{chap:periodic-recoil}

电子与光场的交换已经按模式和有限时间通道建立。若相互作用势还在空间中周期重复，它的空间 Fourier 分量也会限制电子能够获得的动量。例如，一个波矢为 $\vb{G}$ 的 Fourier 分量含有因子 $\ee^{\ii\vb{G}\cdot\vb{r}}$；它乘在入射平面波 $\ee^{\ii\vb{p}\cdot\vb{r}/\hbar}$ 上，就产生动量为 $\vb{p}+\hbar\vb{G}$ 的平面波。周期势引起的动量通道因此可以直接从 Fourier 展开读出。

通道之间的跃迁还受电子能量影响。不同动量对应不同的 $E(\vb{p})$，传播时便积累不同相位。作用时间较短时，许多通道可以共同参与，电子主要获得一个随位置变化的相位；作用时间较长时，不同能量的影响逐渐分开，只有满足共振条件的少数通道能够有效交换振幅。这将把相位光栅与 Bragg 衍射联系起来。\cite{Huang2023Recoil,KarnieliFan2023JC,Ruimy2025Perspective}

我们先用空间平移对称性整理全部动量通道，求出每组通道的演化，再讨论短时和少态极限。若哈密顿量还随时间周期重复，则进一步使用 Floquet 理论。空间周期与时间周期各自产生什么量子数，需要在这两个步骤中分别说明。

\section{周期核与 Bloch 纤维}

对周期势作 Fourier 展开后，只有一组离散波矢参与动量交换。于是，从某个初始动量出发，只能到达与它相差这些离散步长的动量态。我们可以把相互连通的态分成一组，再对每一组独立求解。固体物理中的 Bloch 分解就是这个分组过程；本节从连续动量方程推导它，并说明准动量怎样给各组编号。

\subsection{连续动量一般动力学方程}

在写周期势以前，先保留任意势能算符的连续动量矩阵元。这样可以明确看出，哪些动量联系由一般 Schr\"odinger 方程给出，哪些限制来自周期对称性。

先把晶格语言与真实动量对应起来。\emph{直接晶格}是所有离散平移矢量的集合；若把三个原胞基矢记为 $\vb{a}_i$，则
\begin{equation*}
\mathcal L=\left\{\vb{R}=\sum_{i=1}^3n_i\vb{a}_i\mid n_i\in\ZZ\right\},
\end{equation*}
其中 $\vb{a}_i$ 具有长度量纲。与之配对的\emph{倒格}定义在 Fourier 波矢空间，由所有使 $\ee^{\ii\vb{G}\cdot\vb{R}}=1$ 对每个 $\vb{R}\in\mathcal L$ 成立的离散波矢集合。倒格基由
\begin{equation*}
\vb{a}_i\cdot\vb{b}_j=2\pi\delta_{ij}
\end{equation*}
定义，并生成 $\mathcal L^*$。因此 $\vb{G}\in\mathcal L^*$ 具有逆长度量纲，电子实际交换的机械动量是 $\hbar\vb{G}$。第一布里渊区在波数空间记为 $\mathcal B_{\kappa}$，在机械动量空间记为
\begin{equation*}
\mathcal B_p\equiv\hbar\mathcal B_{\kappa}.
\end{equation*}
其中 $\vb{p}_{\rm B}\in\mathcal B_p$ 是\emph{准动量}代表元，而 $\vb{p}_{\rm B}+\hbar\vb{G}$ 是同一纤维中的真实平面波动量。准动量是所有相差一个倒格动量 $\hbar\vb{G}$ 的平面波等价类在第一布里渊区中的代表元；真实机械动量仍是 $\vb{p}_{\rm B}+\hbar\vb{G}$。第一布里渊区因此只是倒格周期下的一块基本胞，用来避免重复标记同一纤维。

一般投影后的电子动力学若仍含显著环境记忆，应继续使用\chapref{chap:general-interaction}的能量依赖自能或时间卷积核；周期对称性本身并不会消除这种记忆。只有当记忆控制量允许把环境作用化成时间局域单体算符时，周期问题才可写成相应的哈密顿量形式。设电子已经受控地限制在单电子正能扇区，并写
\begin{equation}
{
\hat H(t)=\hat H_e+\hat V(t),
\qquad
\hat H_e|\vb{p},s\rangle=E(\vb{p})|\vb{p},s\rangle .}
\label{eq:arbitrary_single_electron_hamiltonian}
\end{equation}
$\hat V(t)$ 是该扇区中的厄米单体算符；若它来自消去其他量子自由度，则只有在记忆、自能量依赖和环境反作用已经可忽略时才可把精确有效作用化成这种瞬时形式。展开
\begin{equation*}
|\psi(t)\rangle=\sum_s\int\dd^3 p\,c_s(\vb{p},t)|\vb{p},s\rangle,
\end{equation*}
并定义
\begin{equation*}
V_{ss'}(\vb{p},\vb{p}';t)
\equiv\langle\vb{p},s|\hat V(t)|\vb{p}',s'\rangle .
\end{equation*}
连续动量振幅严格满足
\begin{equation}
{
\ii\hbar\partial_t c_s(\vb{p},t)
=E(\vb{p})c_s(\vb{p},t)
+\sum_{s'}\int\dd^3 p'\,
V_{ss'}(\vb{p},\vb{p}';t)c_{s'}(\vb{p}',t).}
\label{eq:arbitrary_potential_continuous_momentum_equation}
\end{equation}
右边第一项给每个动量分量积累自由相位；第二项把所有 $\vb{p}'$ 分量按矩阵元 $V_{ss'}$ 加权后送到 $\vb{p}$。当势具有周期性时，第二项中的积分将只连接特定的动量差，这正是下面进行分组的依据。

若 $\hat V(t)$ 是自旋无关局域标量势 $U(\vb{r},t)$，则矩阵元只依赖动量差。以 $\vb{\kappa}$ 表示傅里叶波矢，$[\vb{\kappa}]={\rm m^{-1}}$，定义
\begin{equation*}
\widetilde U(\vb{\kappa},t)=\int\dd^3 r\,U(\vb{r},t)\ee^{-\ii\vb{\kappa}\cdot\vb{r}},
\end{equation*}
有
\begin{equation*}
V_{ss'}(\vb{p},\vb{p}';t)
=\delta_{ss'}\frac{1}{(2\pi\hbar)^3}
\widetilde U\!\left(\frac{\vb{p}-\vb{p}'}{\hbar},t\right).
\end{equation*}
因此局域周期势的倒格选择定则可以直接读成傅里叶谱支撑；但同一个选择定则对一般非局域周期单体算符仍由平移对称性成立，而不依赖局域势表示。

现在加入离散平移对称性 $[\hat V(t),\hat T_{\vb{R}}]=0$。平面波是平移算符本征态，所以非零矩阵元必须满足
\begin{equation}
{
\vb{p}-\vb{p}'=\hbar\vb{G},
\qquad
\vb{G}\in\mathcal L^*.}
\label{eq:periodic_momentum_selection_rule}
\end{equation}

这个选择定则最容易在一维中建立直觉。若周期为 $a$，则直接晶格为 $R=na$，倒格波数为 $G_m=2\pi m/a$。因此一个周期势不会把入射动量 $p$ 连到任意 $p'$，而只能连接
\[
 p'=p-m\frac{2\pi\hbar}{a},\qquad m\in\ZZ.
\]
离散平移对称性把允许的机械动量转移量子化为 $2\pi\hbar/a$；三维情形则由倒格矢 $\vb{G}$ 取代整数 $m$。周期性因此先生成整条倒格动量格，真实色散与有限作用时间再决定其中哪些通道能够谱隔离。Bragg 两态正是在这一完整动量格上进一步得到的近共振子空间。
这就是准动量模倒格矢守恒。它只说明连续 Hilbert 空间分解成互不耦合的纤维，并不意味着每个纤维只有两个或几个通道；通道截断还需要额外的谱隔离与误差控制。

\begin{proposition}[晶格平移不变算符的动量选择规则]
设单体算符在共同域上与每个晶格平移对易。其动量核的分布支撑只允许 $\vb{p}-\vb{p}\prime=\hbar\vb{G}$，$\vb{G}\in\mathcal L^*$，即\eqnrefs{eq:periodic_momentum_selection_rule}。局域周期势是此一般规则的 Fourier 实现；时间依赖不破坏瞬时的空间规则。
\end{proposition}
\begin{proof}
在完整动量--自旋基上展开
\begin{equation*}
|\psi(t)\rangle=\sum_{s'}\int\dd^3 p'\,
c_{s'}(\vb{p}',t)|\vb{p}',s'\rangle .
\end{equation*}
代入 Schr\"odinger 方程并左乘 $\langle\vb{p},s|$。自由项由
\begin{equation*}
\langle\vb{p},s|\hat H_e|\vb{p}',s'\rangle
=E(\vb{p})\delta_{ss'}\delta^{(3)}(\vb{p}-\vb{p}')
\end{equation*}
化为 $E(\vb{p})c_s(\vb{p},t)$，而相互作用项为
\begin{align*}
\langle\vb{p},s|\hat V(t)|\psi(t)\rangle
&=\sum_{s'}\int\dd^3 p'\,
\langle\vb{p},s|\hat V(t)|\vb{p}',s'\rangle c_{s'}(\vb{p}',t)\\
&=\sum_{s'}\int\dd^3 p'\,
V_{ss'}(\vb{p},\vb{p}';t)c_{s'}(\vb{p}',t),
\end{align*}
两项合并即得正文\eqnrefs{eq:arbitrary_potential_continuous_momentum_equation}。

再施加 $\hat T_{\vb{R}}^\dagger\hat V\hat T_{\vb{R}}=\hat V$。平移动量本征态满足
\begin{equation*}
\hat T_{\vb{R}}|\vb{p},s\rangle
=\ee^{-\ii\vb{p}\cdot\vb{R}/\hbar}|\vb{p},s\rangle,
\qquad
\langle\vb{p},s|\hat T_{\vb{R}}^\dagger
=\ee^{+\ii\vb{p}\cdot\vb{R}/\hbar}
\langle\vb{p},s|.
\end{equation*}
把单位算符形式的平移对称性插入矩阵元：
\begin{align*}
V_{ss'}(\vb{p},\vb{p}';t)
&=\langle\vb{p},s|\hat V(t)|\vb{p}',s'\rangle\\
&=\langle\vb{p},s|\hat T_{\vb{R}}^\dagger
\hat V(t)\hat T_{\vb{R}}|\vb{p}',s'\rangle\\
&=\ee^{\ii(\vb{p}-\vb{p}')\cdot\vb{R}/\hbar}
V_{ss'}(\vb{p},\vb{p}';t).
\end{align*}
移项后得到
\begin{equation*}
\left[
1-\ee^{\ii(\vb{p}-\vb{p}')\cdot\vb{R}/\hbar}
\right]
V_{ss'}(\vb{p},\vb{p}';t)=0
\qquad
\forall\,\vb{R}\in\mathcal L.
\end{equation*}
因此，只要某个矩阵元非零，其相位因子就必须对\emph{每一个}晶格平移都等于一：
\begin{equation*}
\ee^{\ii(\vb{p}-\vb{p}')\cdot\vb{R}/\hbar}=1
\qquad
\forall\,\vb{R}\in\mathcal L.
\end{equation*}
倒格的定义正是
\begin{equation*}
\mathcal L^*
=\left\{\vb{G}\;\middle|\;
\ee^{\ii\vb{G}\cdot\vb{R}}=1
\ \forall\vb{R}\in\mathcal L\right\},
\end{equation*}
所以
\begin{equation*}
\frac{\vb{p}-\vb{p}'}{\hbar}=\vb{G}\in\mathcal L^*,
\end{equation*}
即得到正文\eqnrefs{eq:periodic_momentum_selection_rule}。

局域周期势给出同一结论的 Fourier 表示。若 $U(\vb{r}+\vb{R},t)=U(\vb{r},t)$，其 Fourier 谱只能位于倒格点：
\begin{equation*}
\widetilde U(\vb{\kappa},t)
=\sum_{\vb{G}\in\mathcal L^*}
U_{\vb{G}}(t)(2\pi)^3
\delta^{(3)}(\vb{\kappa}-\vb{G}).
\end{equation*}
代入正文已经定义的局域势矩阵元
\begin{equation*}
V_{ss'}(\vb{p},\vb{p}';t)
=\delta_{ss'}\frac{1}{(2\pi\hbar)^3}
\widetilde U\!\left(\frac{\vb{p}-\vb{p}'}{\hbar},t\right)
\end{equation*}
以后，非零支撑再次要求
$\vb{p}-\vb{p}'=\hbar\vb{G}$。因此“Fourier 谱只含倒格矢”是局域势的具体实现，而平移算符证明说明这个选择定则对一般非局域周期单体算符仍然成立。
\end{proof}


\subsection{Bloch--Floquet 纤维}

“纤维”表示一组只能彼此相差倒格动量的通道。以一维周期 $a$ 为例，令 $G_1=2\pi/a$：$0.3\hbar G_1$、$1.3\hbar G_1$ 和 $2.3\hbar G_1$ 属于同组，$0.4\hbar G_1$ 属于另一组。每组保留一个连续余数，再用整数记录加了多少个倒格动量，便能重新编排整个动量轴。用 $\vb{p}_{\rm B}$ 表示第一布里渊区内的准动量代表，区别于入射束的物理中心动量 $\vb{p}_0$；后者可以远在第一布里渊区之外。

不同纤维“不耦合”只描述哈密顿量的分块结构，并不自动把初态变成纤维间的非相干混合。初态跨越多个准动量并保留相干时，重建空间波函数仍须叠加各纤维振幅。

连续平面波的 $\delta$ 归一化要求把周期哈密顿量写成 Bloch--Floquet 直积分，而不能把单个 $\delta$ 归一化平面波直接当作离散归一态。令 $\mathcal B_p$ 表示第一布里渊区在
de Broglie 动量空间中的基本胞，即 $\mathcal B_p=\hbar\mathcal B_{\kappa}$。任意动量除去边界零测集后都可唯一写成
\begin{equation*}
\vb{p}=\vb{p}_{\rm B}+\hbar\vb{G},
\qquad
\vb{p}_{\rm B}\in\mathcal B_p,
\qquad
\vb{G}\in\mathcal L^*.
\end{equation*}
因此对任意正规化动量波函数 $\psi_s(\vb{p})$，定义纤维变换
\begin{equation*}
(\mathcal U_{\rm BF}\psi)_{\vb{G}s}(\vb{p}_{\rm B})
\equiv
\psi_s(\vb{p}_{\rm B}+\hbar\vb{G}).
\end{equation*}
利用互不重叠（除边界零测集外）的分区
\[
\RR^3
=\bigsqcup_{\vb{G}\in\mathcal L^*}
(\mathcal B_p+\hbar\vb{G}),
\]
先把全空间积分拆成各个平移胞的积分：
\begin{align*}
\sum_s\int_{\RR^3}\dd^3 p\,|\psi_s(\vb{p})|^2
&=\sum_{\vb{G},s}
\int_{\mathcal B_p+\hbar\vb{G}}
\dd^3 p\,|\psi_s(\vb{p})|^2\\
&=\sum_{\vb{G},s}
\int_{\mathcal B_p}\dd^3 p_{\rm B}\,
|\psi_s(\vb{p}_{\rm B}+\hbar\vb{G})|^2\\
&=\int_{\mathcal B_p}\dd^3 p_{\rm B}
\sum_{\vb{G},s}
|\psi_s(\vb{p}_{\rm B}+\hbar\vb{G})|^2\\
&=\int_{\mathcal B_p}\dd^3 p_{\rm B}\,
\| (\mathcal U_{\rm BF}\psi)(\vb{p}_{\rm B})\|_{\ell^2(\mathcal L^*)\otimes\CC^2}^2.
\end{align*}
所以
\begin{equation}
{
L^2(\RR^3,\dd^3 p)\otimes\CC^2
\cong
\int_{\mathcal B_p}^{\oplus}\dd^3 p_{\rm B}\,
\left[\ell^2(\mathcal L^*)\otimes\CC^2\right].}
\label{eq:bloch-floquet-direct-integral-dp68}
\end{equation}
符号 $\cong$ 表示由 $\mathcal U_{\rm BF}$ 实现的幺正同构。这个直积分公式在数学上把一个连续动量波函数重新编排：$\vb{p}_{\rm B}$ 仍是第一布里渊区中的连续变量，而对每一个固定 $\vb{p}_{\rm B}$，所有 $\vb{p}_{\rm B}+\hbar\vb{G}$ 组成一个离散序列。因而一个正规化态不是“选定一个 $\vb{p}_{\rm B}$ 的平面波”，而是一个随 $\vb{p}_{\rm B}$ 变化的 $\ell^2$ 序列族。也正因为如此，固定纤维中的离散矩阵元与纤维之间的 Dirac $\delta$ 必须分开保存；把它们混在一起才会产生形式上的 $\delta(0)$。

平移选择定则保证周期哈密顿量不混合不同 $\vb{p}_{\rm B}$ 纤维，因而对几乎处处的 $\vb{p}_{\rm B}$ 都独立满足
\begin{equation}
{
\ii\hbar\dot c_{\vb{G}s}(\vb{p}_{\rm B},t)
=E(\vb{p}_{\rm B}+\hbar\vb{G})c_{\vb{G}s}(\vb{p}_{\rm B},t)
+\sum_{\vb{G}',s'}
V_{\vb{G}s,\vb{G}'s'}(\vb{p}_{\rm B},t)
 c_{\vb{G}'s'}(\vb{p}_{\rm B},t),}
\label{eq:general_reciprocal_lattice_momentum_equation}
\end{equation}
其中 $V_{\vb{G}s,\vb{G}'s'}(\vb{p}_{\rm B},t)$ 是直积分分解后的纤维矩阵元。连续广义平面波核与它的关系应写成
\begin{equation*}
\langle\vb{p}_{\rm B}+\hbar\vb{G},s|\hat V(t)|
\vb{p}_{\rm B}'+\hbar\vb{G}',s'\rangle
=
\delta^{(3)}(\vb{p}_{\rm B}-\vb{p}_{\rm B}')
V_{\vb{G}s,\vb{G}'s'}(\vb{p}_{\rm B},t),
\end{equation*}
右侧 $\delta$ 因子保留不同 Bloch--Floquet 纤维之间的连续正交性。若采用有限 Born--von Karman 盒，$\vb{p}_{\rm B}$ 离散化，$\delta$ 分布被 Kronecker $\delta$ 取代，$V_{\vb{G}s,\vb{G}'s'}$ 就等于普通归一化盒态的矩阵元；取热力学极限后恢复上述直积分。于是“固定准动量扇区”表示直积分中的一个纤维，物理正规化态仍由跨纤维的可积包络构成。

\begin{theorem}[动量基本胞的 Bloch 直积分]
将动量空间按倒格基本胞划分，则 $(\mathcal U_{\rm BF}\psi)_{\vb{G}s}(\vb{p}_B)=\psi_s(\vb{p}_B+\hbar\vb{G})$ 是 $L^2$ 空间到相应 $\ell^2$ 纤维直积分的幺正等价。周期自伴哈密顿量在此表示中逐准动量解耦，给\eqnrefs{eq:general_reciprocal_lattice_momentum_equation}；基本胞边界仅是零测集。
\end{theorem}
\begin{proof}
考虑无限周期系统，或先取周期边界条件再作热力学极限。势满足离散平移对称性，准动量以第一布里渊区中的代表元标记。定义

\begin{equation*}
(\mathcal U_{\rm BF}\psi)_{\vb{G}s}(\vb{p}_{\rm B})
=\psi_s(\vb{p}_{\rm B}+\hbar\vb{G}).
\end{equation*}
第一步证明它保持内积，而不只保持范数。对任意 $\phi,\psi\in L^2(\RR^3)\otimes\CC^2$，利用倒格基本胞分区，
\begin{align*}
\langle\phi|\psi\rangle
&=\sum_s\int_{\RR^3}\dd^3 p\,
\phi_s^*(\vb{p})\psi_s(\vb{p})\\
&=\sum_{\vb{G},s}
\int_{\mathcal B_p+\hbar\vb{G}}\dd^3 p\,
\phi_s^*(\vb{p})\psi_s(\vb{p}).
\end{align*}
在每个胞作平移变量
$\vb{p}=\vb{p}_{\rm B}+\hbar\vb{G}$。这是雅可比行列式为 $1$ 的平移，所以
\begin{align*}
\langle\phi|\psi\rangle
&=\int_{\mathcal B_p}\dd^3 p_{\rm B}
\sum_{\vb{G},s}
\phi_s^*(\vb{p}_{\rm B}+\hbar\vb{G})
\psi_s(\vb{p}_{\rm B}+\hbar\vb{G})\\
&=\int_{\mathcal B_p}\dd^3 p_{\rm B}\,
\left\langle
(\mathcal U_{\rm BF}\phi)(\vb{p}_{\rm B}),
(\mathcal U_{\rm BF}\psi)(\vb{p}_{\rm B})
\right\rangle_{\ell^2\otimes\CC^2}.
\end{align*}
因此 $\mathcal U_{\rm BF}$ 是等距映射。再写逆变换。除布里渊区边界这一零测集外，每个 $\vb{p}$ 唯一分解为
$\vb{p}=\vb{p}_{\rm B}+\hbar\vb{G}$，故给定纤维序列族 $f_{\vb{G}s}(\vb{p}_{\rm B})$，定义
\begin{equation*}
(\mathcal U_{\rm BF}^{-1}f)_s(\vb{p}_{\rm B}+\hbar\vb{G})
=f_{\vb{G}s}(\vb{p}_{\rm B}).
\end{equation*}
它与正变换逐点互逆，并由内积恒等式保持 $L^2$ 范数，因此直积分分解是幺正等价，而不是形式上的指标重排。最后证明周期哈密顿量的纤维解耦。平移选择定则给
\begin{equation*}
\langle\vb{p},s|V|\vb{p}',s'\rangle\ne0
\quad\Longrightarrow\quad
\vb{p}-\vb{p}'=\hbar\vb{G}_{\rm rec},
\qquad \vb{G}_{\rm rec}\in\mathcal L^*.
\end{equation*}
分别写
$\vb{p}=\vb{p}_{\rm B}+\hbar\vb{G}$ 与
$\vb{p}'=\vb{p}_{\rm B}'+\hbar\vb{G}'$，则
\begin{equation*}
\vb{p}_{\rm B}-\vb{p}_{\rm B}'
=\hbar(\vb{G}_{\rm rec}-\vb{G}+\vb{G}').
\end{equation*}
左边是两个第一布里渊区代表元之差，右边是倒格矢。除边界等价点外，基本胞中唯一可能是
$\vb{p}_{\rm B}=\vb{p}_{\rm B}'$。因此变换后的核必须具有
\begin{equation*}
\delta^{(3)}(\vb{p}_{\rm B}-\vb{p}_{\rm B}')
V_{\vb{G}s,\vb{G}'s'}(\vb{p}_{\rm B}),
\end{equation*}
从而不同准动量纤维互不混合。Born--von Karman 有限盒只把这个连续 Dirac $\delta$ 正规化为离散 Kronecker $\delta$；热力学极限恢复上述直积分，并不存在需要赋值的 $\delta(0)$。
\end{proof}


\subsection{纤维传播子与截断}

固定准动量 $\vb{p}_{\rm B}$ 后，不同组之间已不再耦合，但同一组内仍可能有无限多个倒格通道。若初态只占据其中少数通道，有限时间内的演化可能集中在一个有限窗口。我们先写出这组通道的完整传播子，再估计截去窗口外通道造成的误差。固定 $\vb{p}_{\rm B}$，把态写成
\begin{equation*}
\mathcal H_{\rm lat}=\ell^2(\mathcal L^*)\otimes\CC^2,
\qquad
\vb{c}(t)=\{c_{\vb{G}s}(\vb{p}_{\rm B},t)\}_{\vb{G},s}.
\end{equation*}
把自由色散和周期耦合分别记为 $D$ 与 $V(t)$，则连续一般动力学方程等价于
\begin{equation}
{
\ii\hbar\frac{\dd}{\dd t}\vb{c}(t)
=[D+V(t)]\vb{c}(t),
\qquad
\vb{c}(t)=U(t,t_0)\vb{c}(t_0).}
\label{eq:periodic_lattice_exact_operator_evolution}
\end{equation}

数值计算必须把无限倒格截到有限工作窗。若 $P_L$ 是所选谱投影，真正需要监控的是整个目标时间窗内的外部占据
\begin{equation*}
\epsilon_{\rm tail}(t)
=\|(1-P_L)U(t,t_0)P_L\vb{c}(t_0)\|^2,
\end{equation*}
而不是只检查最终时刻。所有有限维算符范数估计都应理解为在这个已经验证收敛的工作子空间上进行。

在自由作用绘景中，精确传播子满足 Volterra 方程
\begin{equation}
{
U_I(t,t_0)
=I-\frac{\ii}{\hbar}\int_{t_0}^{t}\!\dd t_1\,
V_I(t_1)U_I(t_1,t_0).}
\label{eq:periodic_lattice_volterra_propagator}
\end{equation}
若定义无量纲积分作用量
\begin{equation*}
\Lambda_V(t)=\frac1\hbar\int_{t_0}^{t}\!\dd\tau\,\|V_I(\tau)\|,
\end{equation*}
则 Dyson 级数截到 $N$ 阶的统一余项界为
\begin{equation}
{
\left\|U_I-\sum_{n=0}^{N}U_I^{(n)}\right\|
\le \ee^{\Lambda_V}\frac{\Lambda_V^{N+1}}{(N+1)!}.}
\label{eq:periodic_lattice_dyson_remainder_bound}
\end{equation}
因此“相互作用可以截到低阶”的可计算含义是积分后的 $\Lambda_V$ 足够小，而不是某个瞬时矩阵元看起来小。冻结自由色散、截断时间有序展开和保留长时间相干周期演化，都应从同一个传播子及其误差控制量出发，而不是作为彼此无关的模型。

\begin{theorem}[时间有序 Dyson 级数的范数余项]
设相互作用绘景 $V_I(t)$ 有界且范数可积，$\Lambda_V=\hbar^{-1}\int\|V_I(t)\|\dd t$。第 $n$ 阶时间有序项范数至多 $\Lambda_V^n/n!$，完整级数按算符范数收敛；$N$ 阶截断余项满足 $\|R_N\|\le \ee^{\Lambda_V}\Lambda_V^{N+1}/(N+1)!$，即\eqnrefs{eq:periodic_lattice_dyson_remainder_bound}。
\end{theorem}
\begin{proof}
固定一个 Bloch--Floquet 纤维和实际数值工作窗，把哈密顿量写成
\[
H(t)=D+V(t),
\]
其中 $D$ 是时间无关的自由色散对角算符，$V(t)$ 是周期耦合。自由传播子为
\[
U_D(t,t_0)=\exp\!\left[-\frac{\ii}{\hbar}D(t-t_0)\right].
\]
将完整传播子分解为
\[
U(t,t_0)=U_D(t,t_0)U_I(t,t_0).
\]
对时间求导并代入
$\ii\hbar\partial_tU=(D+V)U$：
\begin{align*}
\ii\hbar\partial_t(U_DU_I)
&=\ii\hbar(\partial_tU_D)U_I+\ii\hbar U_D\partial_tU_I\\
&=DU_DU_I+\ii\hbar U_D\partial_tU_I\\
&=(D+V)U_DU_I.
\end{align*}
消去两侧相同的 $DU_DU_I$，再左乘 $U_D^\dagger$，得到相互作用绘景方程
\[
\ii\hbar\partial_tU_I(t,t_0)=V_I(t)U_I(t,t_0),
\qquad
V_I(t)=U_D^\dagger(t,t_0)V(t)U_D(t,t_0).
\]
由初值 $U_I(t_0,t_0)=I$ 从 $t_0$ 积分到 $t$：
\[
U_I(t,t_0)
=I-\frac{\ii}{\hbar}\int_{t_0}^{t}\dd t_1\,
V_I(t_1)U_I(t_1,t_0),
\]
即正文的 Volterra 方程。把右端的 $U_I(t_1,t_0)$ 再用同一积分式代入一次，得到
\begin{align*}
U_I(t,t_0)
={}&I-\frac{\ii}{\hbar}\int_{t_0}^{t}\dd t_1\,V_I(t_1)\\
&+\left(-\frac{\ii}{\hbar}\right)^2
\int_{t_0}^{t}\dd t_1\int_{t_0}^{t_1}\dd t_2\,
V_I(t_1)V_I(t_2)+\cdots .
\end{align*}
继续迭代后，第 $n$ 阶项为
\[
U_I^{(n)}
=\left(-\frac{\ii}{\hbar}\right)^n
\int_{t_0<t_n<\cdots<t_1<t}
\dd t_1\cdots\dd t_n\,
V_I(t_1)\cdots V_I(t_n).
\]
定义累计无量纲耦合范数
\[
\Lambda_V(t)
\equiv\frac1\hbar\int_{t_0}^{t}\dd s\,\|V_I(s)\|.
\]
由算符范数的次乘性，
\begin{align*}
\|U_I^{(n)}\|
&\le\frac1{\hbar^n}
\int_{t_0<t_n<\cdots<t_1<t}
\dd t_1\cdots\dd t_n
\prod_{j=1}^n\|V_I(t_j)\|.
\end{align*}
被积函数在变量置换下对称，因此有序单纯形恰占完整超立方体积分的 $1/n!$：
\begin{align*}
\|U_I^{(n)}\|
&\le\frac1{n!}
\left[\frac1\hbar\int_{t_0}^{t}\dd s\,\|V_I(s)\|\right]^n\\
&=\frac{\Lambda_V^n}{n!}.
\end{align*}
截到 $N$ 阶的余项 $R_N=U_I-\sum_{n=0}^{N}U_I^{(n)}$ 因而满足
\[
\|R_N\|
\le\sum_{n=N+1}^{\infty}\frac{\Lambda_V^n}{n!}.
\]
令 $n=N+1+m$，并使用
$(N+1+m)!\ge(N+1)!m!$，则
\begin{align*}
\sum_{n=N+1}^{\infty}\frac{\Lambda_V^n}{n!}
&\le\frac{\Lambda_V^{N+1}}{(N+1)!}
\sum_{m=0}^{\infty}\frac{\Lambda_V^m}{m!}\\
&=\ee^{\Lambda_V}\frac{\Lambda_V^{N+1}}{(N+1)!},
\end{align*}
得到正文\eqnrefs{eq:periodic_lattice_dyson_remainder_bound}。
\end{proof}


\section{周期驱动与时间尺度}

空间周期性只决定允许的动量连接；实际动力学还取决于作用时间是否足以分辨不同通道的自由相位。若作用时间远短于相邻已占据通道积累 $O(1)$ 相位差的时间，自由色散在作用窗内近似冻结；若作用时间足够长，谱失谐会选择性抑制离共振通道。常用的短脉冲和长时间极限正是这两个尺度关系的结果。

\subsection{短时相位光栅}

若脉冲足够短，不同动量分量在作用期间积累的自由相位差很小，求解时便可暂时忽略这些相位差。周期势仍能积累有限相位，所以短脉冲不要求调制幅度也很小。把哈密顿量分为自由色散与周期耦合两项，并在所用能量窗口中选取参考能量 $E_*$，定义
\begin{equation}
{
\hat H_{\mathcal W}(t)=D_\perp+V(t),
\qquad
D_\perp\equiv D-E_*\mathbb I.}
\label{eq:periodic_window_hamiltonian_decomposition}
\end{equation}
去掉公共相位以后，完整传播子 $U$ 由 $D_\perp+V$ 生成，而冻结自由传播的 $U_V$ 只由 $V$ 生成。Duhamel 恒等式给出
\begin{equation}
{
\|U(t,t_0)-U_V(t,t_0)\|
\le\frac{|t-t_0|}{\hbar}\|D_\perp\|
\equiv\epsilon_{\rm kin}.}
\label{eq:raman_nath_duhamel_bound}
\end{equation}
因此 Raman--Nath（薄光栅）极限的定义是：相互作用时间短到各个实际占据动量通道之间尚未来得及积累可分辨的自由传播相位。它控制的是作用时间乘以实际访问谱宽，而不是势能幅值。强脉冲仍可产生大的相位和很多衍射级，只要这些级之间的自由相位差在脉冲期间仍小。

\begin{proposition}[有限谱窗内去掉自由色散的误差]
在有限维自伴工作窗内，比较 $D+V(t)$ 与只含 $V(t)$ 的传播，先去掉公共相位 $E_*$。有 $\|U-U_V\|\le T\|D-E_*I\|/\hbar$。选 $E_*=(E_{\max}+E_{\min})/2$ 使界最紧，得到\eqnrefs{eq:raman_nath_duhamel_bound}；窗外传播需另加泄漏误差。
\end{proposition}
\begin{proof}
在实际占据的动量工作窗 $\mathcal W$ 中，去掉任意公共参考能量 $E_*$ 后写
\[
H_{\mathcal W}(t)=D_\perp+V(t),
\qquad
D_\perp=D-E_*I.
\]
令 $U(t,t_0)$ 由 $D_\perp+V(t)$ 生成，$U_V(t,t_0)$ 只由 $V(t)$ 生成：
\begin{align*}
\ii\hbar\partial_tU(t,t_0)&=[D_\perp+V(t)]U(t,t_0),\\
\ii\hbar\partial_tU_V(t,t_0)&=V(t)U_V(t,t_0).
\end{align*}
比较二者时定义
\[
F(s)=U(t,s)U_V(s,t_0).
\]
使用反向传播子的导数
\[
\partial_sU(t,s)=\frac{\ii}{\hbar}U(t,s)[D_\perp+V(s)]
\]
以及
\[
\partial_sU_V(s,t_0)=-\frac{\ii}{\hbar}V(s)U_V(s,t_0),
\]
可得
\begin{align*}
\frac{\dd F}{\dd s}
&=\frac{\ii}{\hbar}U(t,s)[D_\perp+V(s)]U_V(s,t_0)
-\frac{\ii}{\hbar}U(t,s)V(s)U_V(s,t_0)\\
&=\frac{\ii}{\hbar}U(t,s)D_\perp U_V(s,t_0).
\end{align*}
端点满足
\[
F(t)=U_V(t,t_0),
\qquad
F(t_0)=U(t,t_0),
\]
所以积分得到 Duhamel 恒等式
\[
U(t,t_0)-U_V(t,t_0)
=-\frac{\ii}{\hbar}\int_{t_0}^{t}\dd s\,
U(t,s)D_\perp U_V(s,t_0).
\]
正文中的 $H_{\mathcal W}$ 已经表示限制在实际工作窗 $\mathcal W$ 上的有限维算符，因此以下范数都在该工作空间内计算。两个传播子均为酉算符，把 Duhamel 恒等式取算符范数并使用次乘性，得到
\begin{align*}
\|U(t,t_0)-U_V(t,t_0)\|
&\le\frac1\hbar\int_{t_0}^{t}\dd s\,
\|U(t,s)D_\perp U_V(s,t_0)\|\\
&\le\frac1\hbar\int_{t_0}^{t}\dd s\,
\|U(t,s)\|\,\|D_\perp\|\,\|U_V(s,t_0)\|\\
&=\frac{|t-t_0|}{\hbar}\|D_\perp\|.
\end{align*}
若工作窗中的自由能量为 $E_n$，则
\[
\|D_\perp\|=\max_{n\in\mathcal W}|E_n-E_*|.
\]
令 $E_{\max}=\max_{n\in\mathcal W}E_n$、$E_{\min}=\min_{n\in\mathcal W}E_n$。对任意 $E_*$，最大偏差至少满足
\[
\max_{n\in\mathcal W}|E_n-E_*|
\ge \max(|E_{\max}-E_*|,|E_{\min}-E_*|).
\]
右侧在两个端点偏差相等时取得最小值，所以
\[
E_*^{\rm opt}=\frac{E_{\max}+E_{\min}}2,
\qquad
\min_{E_*}\|D_\perp\|
=\frac{E_{\max}-E_{\min}}2.
\]
取 $T=|t-t_0|$，得到
\[
\|U(t,t_0)-U_V(t,t_0)\|
\le\frac{T}{2\hbar}(E_{\max}-E_{\min})
\equiv\epsilon_{\rm kin},
\]
即正文\eqnrefs{eq:raman_nath_duhamel_bound}。
\end{proof}


对一维局域标量周期势
$U(x,t)=\sum_mU_m(t)\ee^{\ii mGx}$，固定一个自旋通道并写 $c_n$ 为动量 $p_0+n\hbar G$ 的振幅，连续动量一般动力学方程化为
\begin{equation}
{
\ii\hbar\dot c_n
=E(p_0+n\hbar G)c_n+\sum_mU_m(t)c_{n-m}.}
\label{eq:general_periodic_momentum_lattice}
\end{equation}
自由色散给出在位能量，势的第 $m$ 个 Fourier 分量给出跨越 $m$ 个倒格格点的跃迁。若显著占据窗口为 $n\in\mathcal W$，最优选择 $E_*$ 为该窗口能量上下界的中点，于是
\[
\epsilon_{\rm kin}
=\frac{T}{2\hbar}
\left[\max_{n\in\mathcal W}E_n-\min_{n\in\mathcal W}E_n\right].
\]
强驱动会扩大 $\mathcal W$，所以“同一个很短的 $T$”也可能随着衍射阶增多而离开薄光栅区。

当 $V(t)$ 是局域标量乘法算符时，不同时刻的 $V(t)$ 彼此对易；冻结自由传播后，时间排序因而完全消失。定义累积相位
\[
\varphi_{\rm g}(x)=\frac1\hbar\int_{t_i}^{t_f}\dd t\,U(x,t),
\]
则
\[
\psi_f(x)=\ee^{-\ii\varphi_{\rm g}(x)}\psi_i(x).
\]
相位光栅的所有衍射振幅就是这个周期函数的 Fourier 系数：
\begin{equation}
{
A_n=\frac1d\int_0^d\dd x\,
\ee^{-\ii nGx}\ee^{-\ii\varphi_{\rm g}(x)},
\qquad d=\frac{2\pi}{G}.}
\label{eq:general_phase_grating_fourier_coefficients}
\end{equation}
若 $\varphi_{\rm g}(x)=\varphi_{\rm g0}\cos Gx$，Jacobi--Anger 恒等式直接给 $A_n=(-\ii)^nJ_n(\varphi_{\rm g0})$。因此 Bessel 定律只在 $\varphi_{\rm g}(x)$ 恰为单一正弦或余弦时出现；一般多谐波光栅仍由\eqnrefs{eq:general_phase_grating_fourier_coefficients} 给出，而不是普遍服从 Bessel 分布。若 $V(t)$ 含自旋矩阵或非局域核，不同时刻可能不对易，此时即使 $D_\perp$ 可以忽略也必须保留时间排序。

\subsection{Floquet--Sambe 空间}

现在考虑持续重复的驱动。设 $H(t+T_d)=H(t)$，其中 $T_d$ 是驱动周期。每经过一个周期，系统都施加同一个传播算符，因此可以先找这个单周期算符的本征态，再描述任意多个周期的演化。它的本征相位将定义\emph{准能量}，一个周期内余下的变化则由周期函数表示，这就是 Floquet 方法。与前面的空间周期问题一样，周期函数可以展开成离散 Fourier 谐波。定义驱动角频率 $\omega_d=2\pi/T_d$，写
\begin{equation*}
H(t)=\sum_{\nu_{\rm F}\in\ZZ}H^{(\nu_{\rm F})}\ee^{+\ii\nu_{\rm F}\omega_d t},
\qquad
|u(t)\rangle=\sum_m|u_m\rangle \ee^{+\ii m\omega_d t}.
\end{equation*}
把波函数的整体相位演化与周期形状分开，写成 $|\Psi(t)\rangle=\ee^{-\ii\varepsilon t/\hbar}|u(t)\rangle$。代回 Schr\"odinger 方程后，时间导数只作用在周期函数 $|u(t)\rangle$ 上，于是原来的含时问题变成周期函数空间中的定态本征问题；这就是 Floquet 构造：
\begin{equation}
{
[H(t)-\ii\hbar\partial_t]|u(t)\rangle
=\varepsilon|u(t)\rangle,
\qquad |u(t+T_d)\rangle=|u(t)\rangle.}
\label{eq:floquet_sambe_eigenproblem}
\end{equation}
周期函数的 Fourier 指标可以像普通离散基指标一样参与矩阵运算。把“原系统 Hilbert 空间 $\mathcal H$”与“整数时间谐波标签 $m\in\ZZ$”做张量积得到的扩展空间 $\mathcal H\otimes\ell^2(\ZZ)$ 通常称为\emph{Sambe 空间}。这里的 $m$ 不是新增的真实粒子数，而只是周期函数的 Fourier 基指标；在这个扩展空间里，原先的含时微分方程被改写成普通矩阵本征值问题。在该基底中，上式对应的矩阵为
\begin{equation}
{
(\mathcal H_F)_{mm'}
=H^{(m-m')}+m\hbar\omega_d\delta_{mm'}.}
\label{eq:general_floquet_sambe_operator_periodic_lattice}
\end{equation}
同一个物理态可以把 $|u(t)\rangle$ 多乘一个整数谐波并同时把 $\varepsilon$ 平移 $\hbar\omega_d$ 的整数倍，所以\emph{准能量} $\varepsilon$ 只在模 $\hbar\omega_d$ 的意义下唯一。它由一周期传播的本征相位定义，刻画“一整个驱动周期之后总共积累多少相位”；这正类似 Bloch 准动量只在倒格矢模意义下唯一。不同 Fourier 指标对应的扩展谱块接近时，就出现周期驱动中的共振。

前面建立的精确投影与受控消元并不会因为周期驱动而改变；唯一变化是被投影的基底现在还要记录时间谐波指标：
\begin{equation*}
|\alpha\rangle
=|\vb{G}_\alpha,m_\alpha,\sigma_\alpha,\ldots\rangle,
\end{equation*}
它同时记录倒格通道、时间谐波和内部量子数。目标投影 $\mathsf P$ 必须在这个完整扩展谱中选取，谱隙 $\Delta_Q$ 也必须在扩展谱中计算，而不能只看裸电子能量。远失谐通道有两种互补的系统消元方式。\emph{Feshbach--Schur 消元}先用投影 $\mathsf P$ 选出希望保留的通道、用 $\mathsf Q=1-\mathsf P$ 收集其余通道，再通过 Schur 补把 $\mathsf Q$ 空间的虚传播精确写成依赖能量的自能；\emph{Schrieffer--Wolff 变换}则在 $\mathsf P$ 与 $\mathsf Q$ 之间存在有限谱隙时，用近恒等的幺正变换逐阶消去两块之间的非对角耦合。前者是能量依赖的精确投影恒等式，后者是有小参数时的阶次展开；两者描述的物理内容都是“离开目标子空间、经过虚中间态、再返回”的修正。将当前扩展谱的 $\mathsf P/\mathsf Q$ 分块代入\chapref{chap:general-interaction}的精确预解式~\eqnrefs{eq:free_electron_feshbach_resolvent_mother}，即可得到能量依赖的有效生成元。只有实际扩展谱隙和非对角耦合允许受控展开时，才进一步取低阶静态结果。Schrieffer--Wolff 的一般幺正分块构造见《高等物理》的投影与有效哈密顿量部分；这里的谱与耦合必须使用当前周期驱动模型。

Floquet 结构特有的新内容是虚路径的选择规则。若第 $a$ 个微观顶角携带时间傅里叶指标 $r_a$ 和倒格矢转移 $\Delta\vb{G}_a$，任意从目标态 $i$ 到 $j$ 的有效路径必须满足
\begin{equation}
{
\sum_a r_a=m_j-m_i,
\qquad
\sum_a\Delta\vb{G}_a=\vb{G}_j-\vb{G}_i.}
\label{eq:floquet-virtual-path-selection-dp19}
\end{equation}
自旋、能带和偏振选择定则进一步把不允许的微观矩阵元置零，因此只剩满足这些选择定则的 Fourier 顶角参与路径求和。当这些 Fourier 指标之和等于 $n$ 时，通常把该有效跃迁简称为“$n$ 光子过程”；这个名称只是在记录净交换的驱动谐波数，它仍由同一周期哈密顿量的多步矩阵元构成。

\begin{lemma}[Sambe 虚路径的谐波与倒格记账]
在固定 $\ee^{+\ii r\omega_dt}$ 的 Fourier 约定下，每个顶角只改变 Sambe 指标 $r$ 及其空间倒格转移。故任意有限阶非零虚路径满足\eqnrefs{eq:floquet-virtual-path-selection-dp19}。能量分母改变路径权重，近共振要求调整目标空间，但不会改变离散选择规则。
\end{lemma}
\begin{proof}
Sambe 基态写成 $|\vb{G},m,\sigma\rangle$。由
\[
H(t)=\sum_r H^{(r)}\ee^{+\ii r\omega_d t}
\]
和正文\eqnrefs{eq:general_floquet_sambe_operator_periodic_lattice}，第 $r$ 个时间 Fourier 分量只连接相差 $r$ 个 Sambe 指标的扇区。若该分量的空间 Fourier 成分再携带倒格矢转移 $\Delta\vb{G}$，对应微观顶角的矩阵元具有结构
\begin{align*}
&\langle \vb{G}',m',\sigma'|
V_{r,\Delta\vb{G}}
|\vb{G},m,\sigma\rangle\\
&\qquad=
\mathcal V^{(r,\Delta\vb{G})}_{\sigma'\sigma}(\vb{G})
\,\delta_{m',m+r}
\,\delta_{\vb{G}',\vb{G}+\Delta\vb{G}}.
\end{align*}
这里 $\mathcal V^{(r,\Delta\vb{G})}_{\sigma'\sigma}$ 包含自旋、能带、偏振等内部矩阵元；若内部选择定则禁止该跃迁，它本身为零。真正负责时间谐波与倒格通道记账的是两个 Kronecker 约束。

先看二阶虚过程。插入一组完整中间态 $|\alpha\rangle=|\vb{G}_\alpha,m_\alpha,\sigma_\alpha\rangle$，典型有效矩阵元为
\[
( H_{\rm eff}^{(2)} )_{ji}
=\sum_{\alpha\in Q}
\frac{V_{j\alpha}V_{\alpha i}}
{E_i-E_\alpha}.
\]
若第一步携带 $(r_1,\Delta\vb{G}_1)$，第二步携带 $(r_2,\Delta\vb{G}_2)$，两个矩阵元分别含
\begin{align*}
V_{\alpha i}&\propto
\delta_{m_\alpha,m_i+r_1}
\delta_{\vb{G}_\alpha,\vb{G}_i+\Delta\vb{G}_1},\\
V_{j\alpha}&\propto
\delta_{m_j,m_\alpha+r_2}
\delta_{\vb{G}_j,\vb{G}_\alpha+\Delta\vb{G}_2}.
\end{align*}
第一个约束先固定中间指标，代入第二个约束便得到
\begin{align*}
m_j-m_i&=r_1+r_2,\\
\vb{G}_j-\vb{G}_i
&=\Delta\vb{G}_1+\Delta\vb{G}_2.
\end{align*}
能量分母只改变这条允许路径的权重，并不改变这两个离散指标约束。

对 $N$ 阶过程，依次插入 $N-1$ 组完备关系。设第 $a$ 个顶角连接 $\alpha_{a-1}\to\alpha_a$，其中 $\alpha_0=i$、$\alpha_N=j$。每个非零矩阵元都要求
\begin{align*}
m_{\alpha_a}-m_{\alpha_{a-1}}&=r_a,\\
\vb{G}_{\alpha_a}-\vb{G}_{\alpha_{a-1}}
&=\Delta\vb{G}_a.
\end{align*}
把 $a=1,\ldots,N$ 的等式逐项相加，中间指标成对消去：
\begin{align*}
\sum_{a=1}^{N}
(m_{\alpha_a}-m_{\alpha_{a-1}})
&=m_j-m_i
=\sum_{a=1}^{N}r_a,\\
\sum_{a=1}^{N}
(\vb{G}_{\alpha_a}-\vb{G}_{\alpha_{a-1}})
&=\vb{G}_j-\vb{G}_i
=\sum_{a=1}^{N}\Delta\vb{G}_a.
\end{align*}
因此任意非零高阶虚路径都满足
\[
{
\sum_a r_a=m_j-m_i,
\qquad
\sum_a\Delta\vb{G}_a=\vb{G}_j-\vb{G}_i,}
\]
即正文\eqnrefs{eq:floquet-virtual-path-selection-dp19}。若某个中间能量分母趋近于零，改变的是 $P/Q$ 子空间划分：该近共振态必须纳入目标空间；上述由 Fourier 矩阵元产生的选择规则本身保持不变。
\end{proof}


在谱隔离条件 $v_{PQ}/\Delta_Q\ll1$ 下，Feshbach/Schrieffer--Wolff 有效哈密顿量给出虚路径振幅：二阶量级为 $v^2/\Delta_Q$，三阶为 $v^3/\Delta_Q^2$，更高阶依次增加一个耦合因子和一个失谐分母。一旦某个中间分母逼近零，该通道必须从 $\mathsf Q$ 空间提升到新的 $\mathsf P$ 空间；共振通道不能继续作为微扰消去项处理。

有限包络或啁啾破坏严格周期性时，没有单一 $\mathcal H_F$ 能替代原问题，应回到时间排序传播子。若仍使用随时间变化的低维子空间，还必须保留投影算符运动产生的非绝热项 $-\ii\hbar U^\dagger\dot U$。因此短脉冲相位光栅、严格周期 Floquet 谱与一般有限包络是同一周期势一般动力学方程的三种不同时间结构，而不是三套彼此独立的模型。

\subsection{高频 Floquet 有效哈密顿量与微运动}

上一小节把周期含时问题改写成不同 Fourier 谐波之间的耦合。若驱动频率很高，不同谐波副本的能量差很大，电子很少真实占据远离目标的副本。然而，先进入这些副本再返回的虚过程仍会改变慢演化。在高频展开中，时间平均 $H^{(0)}$ 给出最低阶哈密顿量，这些往返过程则给出更高阶修正。我们还需要保留一个周期内的快速运动，才能从慢演化恢复任意时刻的态。

为了计算时间平均以外的修正，在前面的 Sambe 矩阵中选取目标子空间为 $m=0$ 谐波扇区，$\mathsf P=\sum_a|a,0\rangle\langle a,0|$；其余 $m\neq0$ 扇区记为 $\mathsf Q$。如果工作窗内所有目标本征能 $E_a$ 与移位副本满足
\begin{equation}
{
\Delta_F
\equiv
\inf_{a,b,\,m\neq0}
|E_a-E_b-m\hbar\omega_d|>0,
\qquad
\epsilon_{\rm Floq}
\equiv
\frac{\sum_{m\neq0}\|H^{(m)}\|}{\Delta_F}\ll1,}
\label{eq:floquet_high_frequency_gap_control}
\end{equation}
则非零 Sambe 扇区可以作受控虚消元。这里真正控制展开的是与所有可能谐波共振的距离 $\Delta_F$，而不是单独要求“$\omega_d$ 很大”；若某个 $E_a-E_b\simeq m\hbar\omega_d$，相应谐波态必须提升到目标 $\mathsf P$ 空间，进入共振 Floquet 问题。

把静态能宽与振荡耦合分开计量。选取公共能量 $E_*$，记
\[
B_F=\|H^{(0)}-E_*I\|,\qquad
A_F=\sum_{m\ne0}\|H^{(m)}\|,\qquad W_F=\hbar\omega_d.
\]
这些量都具有能量量纲。谱隔离足以允许虚消元；若还要把能量分母的首项统一换成 $mW_F$，则需要更强的高频条件 $(B_F+A_F)/W_F\ll1$。按全文采用的 $H(t)=\sum_mH^{(m)}\ee^{+\ii m\omega_dt}$ 约定，得到
\begin{equation}
{
H_{F,\rm eff}
=H^{(0)}
+\sum_{m=1}^{\infty}
\frac{[H^{(m)},H^{(-m)}]}{m\hbar\omega_d}
+O\!\left(\frac{A_F^2B_F+A_F^3}{W_F^2}\right),}
\label{eq:floquet_van_vleck_first_order_our_convention}
\end{equation}
余项常数可依赖固定的谐波数，但不依赖驱动频率。$A_F^2B_F$ 来自虚传播分母中保留静态能量，$A_F^3$ 来自第三个振荡顶角。若不同谐波算符彼此对易，$O(\omega_d^{-1})$ 项消失，单纯时间平均才在这一阶成立；静态部分若不与驱动对易，下一阶仍可能非零。

\begin{proposition}[有限谐波高频展开的首阶与微运动]
设有限维自伴周期哈密顿量只有有限个 Fourier 谐波，且频率大于保留能标与耦合，目标 Sambe 块和补块保持非共振谱隙。按本章正指数 Fourier 约定，首阶有效哈密顿量为\eqnrefs{eq:floquet_van_vleck_first_order_our_convention}，首阶周期微运动生成元为\eqnrefs{eq:floquet_micromotion_kick_first_order}。余项按 $\omega_d^{-2}$ 控制；无限谐波需补相应范数求和条件。
\end{proposition}
\begin{proof}
为隔离高频结构，先把 $m=0$ Sambe 扇区作为目标空间。由正文\eqnrefs{eq:general_floquet_sambe_operator_periodic_lattice}，零阶块为
\[
\mathsf P\mathcal H_F\mathsf P=H^{(0)},
\]
而 Sambe 矩阵的第 $m$ 个对角块严格为
\[
\mathsf Q_m\mathcal H_F\mathsf Q_m
=\mathsf Q_m\bigl(H^{(0)}+m\hbar\omega_d\bigr)\mathsf Q_m.
\]
非零 Fourier 分量 $H^{(n\ne0)}$ 只连接不同 Sambe 扇区，因此不属于这一对角块；它们随后通过 Feshbach 自能进入虚跃迁修正。
从 $m=0$ 到 $m$ 的 Fourier 顶角由矩阵元
\[
(\mathcal H_F)_{m0}=H^{(m)},
\qquad
(\mathcal H_F)_{0m}=H^{(-m)}
\]
给出。精确 Feshbach 自能为
\[
\Sigma_P(E)
=\mathsf P\mathcal H_F\mathsf Q
(E-\mathsf Q\mathcal H_F\mathsf Q)^{-1}
\mathsf Q\mathcal H_F\mathsf P.
\]
在高频非共振条件下，先忽略 $\mathsf Q$ 内部的次级 Fourier 耦合，并把分母按 $m\hbar\omega_d$ 展开。对 $m>0$ 扇区，
\[
(E-H^{(0)}-m\hbar\omega_d)^{-1}
=-\frac1{m\hbar\omega_d}
+O\!\left(\frac{\|E-H^{(0)}\|}{(m\hbar\omega_d)^2}\right),
\]
所以该扇区贡献
\[
-\frac{H^{(-m)}H^{(m)}}{m\hbar\omega_d}.
\]
相应的 $-m$ 扇区分母符号相反，给
\[
+\frac{H^{(m)}H^{(-m)}}{m\hbar\omega_d}.
\]
把一对 $\pm m$ 扇区相加：
\begin{align*}
\Sigma_m^{(2)}
&=\frac{H^{(m)}H^{(-m)}-H^{(-m)}H^{(m)}}
{m\hbar\omega_d}\\
&=\frac{[H^{(m)},H^{(-m)}]}
{m\hbar\omega_d}.
\end{align*}
再对所有 $m>0$ 求和，就得到正文\eqnrefs{eq:floquet_van_vleck_first_order_our_convention} 的首阶。这里的对易子符号与 Fourier 约定绑定；若改用 $\ee^{-\ii m\omega_dt}$，相应写法会交换 $m$ 与 $-m$。

余项有两种来源。对目标低能分支，$|E-E_*|=O(B_F+A_F)$。令 $X=E-H^{(0)}$，当 $\|X\|<|m|W_F$ 时，预解式恒等式给
\[
\left\|(X-mW_F)^{-1}+\frac{I}{mW_F}\right\|
\le\frac{\|X\|}{|m|W_F(|m|W_F-\|X\|)}.
\]
左右各乘一个 Fourier 顶角并求和，这一分母修正为
$O[A_F^2(B_F+A_F)/W_F^2]$。再保留一次 $Q$ 块内部 Fourier 跃迁时，三个顶角与两个高频分母给 $O(A_F^3/W_F^2)$；有限谐波的更长路径由 Neumann 级数控制。归一化目标分量所需的近恒等幺正变换同样从 $O(A_F^2/W_F^2)$ 开始，和静态能标相乘后也计入 $A_F^2B_F/W_F^2$。这就给出正文的余项阶数。若只有谱隔离而静态能宽不能相对 $W_F$ 忽略，应保留原有能量分母；若出现近共振谐波，则把该扇区并入目标空间。

最低阶微运动可从一个周期近恒等变换得到。令
\[
|\Psi(t)\rangle=\ee^{-\ii K(t)}|\Phi(t)\rangle,
\qquad K(t+T_d)=K(t),
\]
并要求 $K=O(\omega_d^{-1})$ 消去变换后哈密顿量的非零 Fourier 一阶项。展开
\[
H'=
\ee^{+\ii K}H\ee^{-\ii K}
-\ii\hbar\ee^{+\ii K}\partial_t\ee^{-\ii K}
=H-\hbar\dot K+O(\omega_d^{-1}H^2).
\]
逐个非零谐波要求
\[
H^{(m)}\ee^{+\ii m\omega_dt}
-\hbar\partial_tK_m(t)=0,
\]
于是
\[
K_m(t)=\frac{H^{(m)}\ee^{+\ii m\omega_dt}}
{\ii m\hbar\omega_d},
\]
求和即得正文\eqnrefs{eq:floquet_micromotion_kick_first_order}。
\end{proof}


周期内部还存在快速运动。在有限工作窗内，选定单周期传播子的厄米对数 $H_F$，使 $U(T_d,0)=\exp(-\ii H_FT_d/\hbar)$。定义周期幺正算符 $P_F(t)=U(t,0)\exp(+\ii H_Ft/\hbar)$。由 $U(t+T_d,0)=U(t,0)U(T_d,0)$，直接得到 $P_F(t+T_d)=P_F(t)$；精确的 Floquet 分解因此为
\begin{equation}
{
U(t,t_0)
=P_F(t)\ee^{-\ii H_F(t-t_0)/\hbar}P_F^\dagger(t_0),
\qquad P_F(t+T_d)=P_F(t),}
\label{eq:floquet_micromotion_factorization}
\end{equation}
对任意常数幺正算符 $C_F$，同时替换 $P_F\mapsto P_FC_F$、$H_F\mapsto C_F^\dagger H_FC_F$ 不改变这一传播子。因而把 $P_F(0)$ 固定为恒等与把微运动的平均生成元固定为零，是不同的表示选择；后面的高频展开使用后者。
周期幺正 $P_F$ 不自动具有连续周期的厄米对数。若选定的微运动表示在整个周期内接近恒等，例如 $\sup_t\|P_F(t)-I\|<1$，其谱不会穿过负实轴，主对数便连续且周期；这时可定义 $K(t)=\ii\log P_F(t)$，写成 $P_F(t)=\exp[-\ii K(t)]$。更一般地，也可在确实存在相容连续周期对数支路时使用 $K$。在这种表示中，精确分解等价于
\[
U(t,t_0)=\ee^{-\ii K(t)}\ee^{-\ii H_F(t-t_0)/\hbar}\ee^{+\ii K(t_0)}.
\]
高频展开通常选择近恒等微运动，再以 $H_{F,\rm eff}$ 近似 $H_F$；其最低阶无量纲微运动生成元为
\begin{equation}
{
K^{(1)}(t)
=\sum_{m\neq0}
\frac{H^{(m)}\ee^{+\ii m\omega_dt}}
{\ii m\hbar\omega_d}.}
\label{eq:floquet_micromotion_kick_first_order}
\end{equation}
在整数周期后，$P_F(t_0+nT_d)=P_F(t_0)$，所以实际单周期生成元是 $P_F(t_0)H_FP_F^\dagger(t_0)$。在近恒等的高频表示中，它近似为 $\ee^{-\ii K(t_0)}H_{F,\rm eff}\ee^{+\ii K(t_0)}$，且一般 $K(t_0)\ne0$。读出慢演化时，初态与观测算符也必须作同一变换；亚周期测量则保留两端不同的微运动。若静态哈密顿量的余项范数为 $r_F$，Duhamel 界给出的长期相位误差还含 $|t-t_0|r_F/\hbar$。频率很高因而不能单独保证任意长时间的传播都准确。

\begin{exbox}[同一旋转驱动中的静态能宽修正]
在有限两态工作窗内，取 $\hbar\omega_d=W_F>0$，并设
\[
H(t)=B\sigma_z+a[\sigma_x\cos(\omega_dt)+\sigma_y\sin(\omega_dt)],
\qquad a,B\in\RR,\quad 0<B<W_F/2.
\]
求接近 $+B$ 的准能量分支，核对首阶高频修正及其下一阶能标。

\textbf{解。}Fourier 顶角为 $H^{(+1)}=a\sigma_-$、$H^{(-1)}=a\sigma_+$，所以首阶有效哈密顿量是 $(B-a^2/W_F)\sigma_z$。这个符号也可直接用精确含时解核对。令 $R(t)=\exp(-\ii\omega_dt\sigma_z/2)$，写 $|\psi\rangle=R|\phi\rangle$，则
\[
H_{\rm rot}=R^\dagger HR-\ii\hbar R^\dagger\dot R
=(B-W_F/2)\sigma_z+a\sigma_x.
\]
因为 $R(T_d)=-I$，一周期相位还包含这个负号。选取连续趋向 $B$ 的准能量支路，得到
\[
q_+=\frac{W_F}{2}-\sqrt{(W_F/2-B)^2+a^2}
=B-\frac{a^2}{W_F-2B}
+O\!\left(\frac{a^4}{(W_F-2B)^3}\right).
\]
再以 $B/W_F$ 展开分母，
\[
q_+=B-\frac{a^2}{W_F}-\frac{2Ba^2}{W_F^2}
+O\!\left(\frac{a^2B^2+a^4}{W_F^3}\right).
\]
因此下一阶确有静态能宽乘两个驱动顶角的 $Ba^2/W_F^2$ 项。当 $a\ll B$ 时，它会大于仅用 $a^3/W_F^2$ 估计的尺度。相反，$W_F\to2B$ 时分母变小，问题转入真实谐波共振，必须保留相应扩展态。
\end{exbox}

\subsection{驻波周期势的参数化与两种时间尺度}

作为一般周期势的应用，考虑两束反向传播的光形成的驻波。由最小耦合哈密顿量求出它的空间 Fourier 分量后，即可代入\eqnrefs{eq:general_periodic_momentum_lattice}。改变脉冲持续时间，会让同一装置从多通道的相位光栅逐渐进入选择性共振的区域；我们用驻波的耦合强度和反冲能量计算这一变化。

\begin{exbox}[Kapitza--Dirac 两种衍射极限]
Kapitza--Dirac 衍射指自由电子被空间周期的光学驻波相干衍射，电子获得的横向机械动量以光场倒格动量的整数倍变化。它提供同一周期动量格在两个时间尺度上的直接例子：短作用时间进入 Raman--Nath 薄光栅区，许多衍射级可以同时出现；长作用时间能够分辨反冲能量差时进入 Bragg 谱选择区，通常只剩一对或少数近共振动量通道。

取电子主要沿 $z$ 传播，驻波沿 $x$ 方向形成，并令偏振沿 $y$。记单束激光的波数为 $\kappa_{\rm L}=\omega/c$，则
\[
A_y(x,t)=A_0\cos(\kappa_{\rm L}x)\cos(\omega t),
\qquad E_0=\omega A_0.
\]
求周期平均后的驻波势，并分别计算短脉冲衍射与谱选择二态区的跃迁概率。

\textbf{解。} 从非相对论最小耦合哈密顿量
\begin{equation}
{
\hat H=\frac{(\hat{\vb{p}}-q_{\mathrm e}\vb{A})^2}{2m_{\mathrm e}}.}
\label{eq:kd_minimal_coupling_start}
\end{equation}
出发，中心几何取 $p_y=0$ 时线性 $\vb{A}\!\cdot\!\vb{p}$ 顶点精确消失。若电子具有有限横向动量宽度 $\Delta p_y$，可用
\begin{equation*}
 \epsilon_{p_y}\equiv\frac{2\Delta p_y}{|q_{\mathrm e}|A_0}\ll1
\end{equation*}
控制线性顶点相对于 $A^2$ 驱动的幅度。再令光场包络变化时间为 $\tau_A$，则光学周期平均至少要求
\begin{equation*}
 \epsilon_{\rm env,KD}=\frac1{\omega\tau_A}\ll1,
 \qquad
 \epsilon_{\rm micro,KD}
 =\max\!\left(\frac{\omega_{\rm r}}{\omega},
 \frac{U_p}{\hbar\omega}\right)\ll1,
\end{equation*}
前者控制包络在单个载频周期内的变化，后者控制电子反冲与有质动力响应在 $2\omega$ 微运动时间尺度上的累积。保留这些条件下 $A^2$ 的周期平均首项，定义有质动力能标
\[
U_p\equiv\frac{q_{\mathrm e}^2E_0^2}{4m_{\mathrm e}\omega^2},
\]
并得到有效驻波势
\begin{equation}
{
U_{\rm KD}(x)=U_p\cos^2(\kappa_{\rm L}x)
=\frac{U_p}{2}\bigl[1+\cos(\kappa_{\rm KD}x)\bigr].}
\label{eq:kd_ponderomotive_lattice}
\end{equation}
常数 $U_p/2$ 只产生全局相位；真正负责衍射的是倒格波矢 $\kappa_{\rm KD}\equiv2\kappa_{\rm L}$ 的 Fourier 分量。因此 KD 的相邻格点相差横向动量 $\hbar\kappa_{\rm KD}$，而不是 PINEM 中的纵向能量 $\hbar\omega$。对亚周期驱动或显式载波相位测量，周期平均不成立，应回到原始含时最小耦合哈密顿量。

短脉冲时，前述 Raman--Nath 条件~\eqnrefs{eq:raman_nath_duhamel_bound} 要求作用时间内\emph{实际被占据}的动量级之间尚未积累可分辨的自由相位。忽略这一窗口内的横向自由传播后，\eqnrefs{eq:general_phase_grating_fourier_coefficients} 直接给出
\begin{equation}
{
\psi_f(x)\doteq
\ee^{-\ii\Theta_{\rm KD}\cos(\kappa_{\rm KD}x)}\psi_i(x),
\qquad
\Theta_{\rm KD}\equiv\frac{U_p\tau}{2\hbar},}
\label{eq:kd_raman_nath_phase}
\end{equation}
其中 $\doteq$ 表示忽略共同相位。由于该光栅只有一个余弦谐波，傅里叶系数退化为 Bessel 函数，故
\begin{equation}
{
\psi_f(x)=\sum_{n\in\ZZ}(-\ii)^nJ_n(\Theta_{\rm KD})
\ee^{\ii n\kappa_{\rm KD}x}\psi_i(x),
\qquad
P_n=J_n^2(\Theta_{\rm KD}).}
\label{eq:kd_raman_nath_bessel}
\end{equation}
若入射横向动量为 $p_x$，第 $n$ 级峰位于 $p_x+n\hbar\kappa_{\rm KD}$。定义单光子反冲频率
\[
\omega_{\rm KD}\equiv\frac{\hbar\kappa_{\rm L}^2}{2m_{\mathrm e}},
\]
则从 $p_x\simeq0$ 出发，第 $n$ 级相对自由能为 $4n^2\hbar\omega_{\rm KD}$。如果脉冲显著占据到 $|n|\le n_{\max}$，一个直接可检验的充分条件是
\begin{equation}
{\epsilon_{\rm RN}^{\rm KD}
\equiv4n_{\max}^2\omega_{\rm KD}\tau\ll1.}
\label{eq:kd_raman_nath_control}
\end{equation}
所以“短脉冲”并非单指 $\tau$ 小，而是指作用时间短到不能解析当前实际访问的反冲谱宽。

作用时间增长后必须保留自由色散。把波函数展开到倒格动量格
\[
\psi(x,t)=\sum_{n\in\ZZ}c_n(t)
\exp\!\left[\frac{\ii}{\hbar}(p_x+n\hbar\kappa_{\rm KD})x\right],
\]
\eqnrefs{eq:general_periodic_momentum_lattice} 对当前单谐波势化为
\begin{equation}
{
\ii\hbar\dot c_n
=\frac{(p_x+n\hbar\kappa_{\rm KD})^2}{2m_{\mathrm e}}c_n
+\frac{U_p}{4}\bigl(c_{n-1}+c_{n+1}\bigr).}
\label{eq:kd_exact_momentum_chain}
\end{equation}
其中在位能量由自由反冲色散给出，最近邻跃迁 $U_p/4$ 来自 $\cos(\kappa_{\rm KD}x)$ 的两个傅里叶分量。长时间并不会自动产生二能级模型；还必须有一对近简并格点，并且它们与其余通道之间存在足够大的谱隙。

一阶 Bragg 工作点取 $p_x=-\hbar\kappa_{\rm L}$，于是 $|-\hbar\kappa_{\rm L}\rangle$ 与 $|+\hbar\kappa_{\rm L}\rangle$ 的自由能同为
\[
E_{\rm KD}\equiv\frac{\hbar^2\kappa_{\rm L}^2}{2m_{\mathrm e}}.
\]
若最近的 $|\pm3\hbar\kappa_{\rm L}\rangle$ 等通道与目标双态之间的最小失谐记为 $\Delta_{\rm off}$，则谱投影需要
\[
\epsilon_{\rm B}^{\rm KD}\equiv\frac{U_p}{4\Delta_{\rm off}}\ll1.
\]
在这一受控极限中，Feshbach/Schrieffer--Wolff 约化给最低阶有效哈密顿量
\begin{equation}
{
H_{\rm KD}^{(2)}
=E_{\rm KD}\mathbb I_2+\frac{U_p}{4}\sigma_x
+O\!\left(\frac{U_p^2}{\Delta_{\rm off}}\right),}
\label{eq:kd_bragg_two_level_hamiltonian}
\end{equation}
因此从 $|-\hbar\kappa_{\rm L}\rangle$ 出发，最低阶 Bragg 转移概率为
\begin{equation}
{
P_{-\hbar\kappa_{\rm L}\rightarrow+\hbar\kappa_{\rm L}}(\tau)
=\sin^2\!\left(\frac{U_p\tau}{4\hbar}\right).}
\label{eq:kd_bragg_rabi_probability}
\end{equation}
在完全对称的一阶工作点，两个最近离共振通道产生相同的二阶光移，因此只改共同相位；离开对称点后，差分光移与高阶有效耦合必须一起保留。

\medskip

\eqnrefs{eq:kd_minimal_coupling_start} 同时包含线性 $A\!\cdot\!p$ 与二次 $A^2$ 项。当前驻波几何在中心轨迹上消去主导线性横向顶角；再对准单色载频做周期平均，$A^2$ 的常数与 $2\kappa_{\rm L}$ 空间谐波保留下来。显式完成这一展开与时间平均，从而得到\eqnrefs{eq:kd_ponderomotive_lattice}；若包络在单周期内显著变化，则最后的周期平均不能使用。展开最小耦合，
\[
\hat H=\frac{\hat{\vb{p}}^{\,2}}{2m_{\mathrm e}}
-\frac{q_{\mathrm e}}{2m_{\mathrm e}}\bigl(\hat{\vb{p}}\!\cdot\!\vb{A}+\vb{A}\!\cdot\!\hat{\vb{p}}\bigr)
+\frac{q_{\mathrm e}^2A^2}{2m_{\mathrm e}}.
\]
当前几何中 $A_y=A_y(x,t)$ 且主电子动量没有 $y$ 分量，因此线性项在理想平面波近似下不驱动主导横向格点交换。对 $A^2$ 做一个光学周期平均，
\[
\langle A_y^2\rangle_t
=A_0^2\cos^2(\kappa_{\rm L}x)\langle\cos^2\omega t\rangle_t
=\frac{A_0^2}{2}\cos^2(\kappa_{\rm L}x).
\]
于是
\[
\frac{q_{\mathrm e}^2}{2m_{\mathrm e}}\langle A_y^2\rangle_t
=\frac{q_{\mathrm e}^2A_0^2}{4m_{\mathrm e}}\cos^2(\kappa_{\rm L}x)
=\frac{q_{\mathrm e}^2E_0^2}{4m_{\mathrm e}\omega^2}\cos^2(\kappa_{\rm L}x),
\]
即得到 $U_p$ 与\eqnrefs{eq:kd_ponderomotive_lattice}。周期平均来自：若包络或测量解析到载频尺度，最后一步不能做。

\medskip

\eqnrefs{eq:kd_raman_nath_phase} 已经把短脉冲驻波作用化成一个余弦相位因子。这个相位光栅只有单一空间谐波，因此 Jacobi--Anger 展开直接给出各倒格动量阶的振幅：
\[
\ee^{-\ii\Theta\cos\phi}
=\sum_{n\in\ZZ}(-\ii)^nJ_n(\Theta)\ee^{\ii n\phi}.
\]
取 $\Theta=\Theta_{\rm KD}$、$\phi=\kappa_{\rm KD}x$，每个指数项把横向动量平移 $n\hbar\kappa_{\rm KD}$，于是第 $n$ 阶振幅为 $(-\ii)^nJ_n(\Theta_{\rm KD})$。取模平方便得到\eqnrefs{eq:kd_raman_nath_bessel}。

\medskip

\eqnrefs{eq:raman_nath_duhamel_bound} 表明，Raman--Nath 近似真正控制的是实际访问动量窗内的自由谱宽乘以脉冲时间。对 Kapitza--Dirac 格点，$p_x\simeq0$ 时第 $n$ 阶相对中心态的自由能差为
\[
\Delta E_n
=\frac{(n\hbar\kappa_{\rm KD})^2}{2m_{\mathrm e}}
=4n^2\hbar\omega_{\rm KD}.
\]
若显著占据只延伸到 $|n|\le n_{\max}$，Duhamel 界中的最大自由相位为
\[
\max_{|n|\le n_{\max}}\frac{|\Delta E_n|\tau}{\hbar}
=4n_{\max}^2\omega_{\rm KD}\tau.
\]
要求这一量远小于一，就得到\eqnrefs{eq:kd_raman_nath_control}。因此“短脉冲”同时依赖 $\tau$ 和实际占据的最高衍射阶，而不是一个脱离动力学带宽的时间条件。

\medskip

离开 Raman--Nath 极限后，各动量阶在作用时间内会积累可分辨的自由相位，因此必须保留真实反冲失谐。精确起点仍是多通道动量链~\eqnrefs{eq:kd_exact_momentum_chain}。

\emph{Bragg 区域}对应一种特殊的谱结构：周期势近似单色，一对动量态的自由能差恰好被一个倒格动量散射过程补偿，因而二者近共振；其余动量阶与这对目标态之间仍存在明显失谐。实验上它表现为两个选定动量通道之间的相干振荡，而不是短脉冲相位光栅那样同时填充许多衍射级。此时可把近共振双态作为慢子空间，用 Feshbach--Schur 或 Schrieffer--Wolff 消去远失谐通道；泄漏误差由前述谱隔离控制量给出。把倒格展开代入
\[
\ii\hbar\partial_t\psi
=\left[-\frac{\hbar^2}{2m_{\mathrm e}}\partial_x^2+
\frac{U_p}{2}\cos(\kappa_{\rm KD}x)\right]\psi,
\]
并使用
\[
\cos(\kappa_{\rm KD}x)=\frac12\bigl(\ee^{\ii\kappa_{\rm KD}x}+\ee^{-\ii\kappa_{\rm KD}x}\bigr),
\]
投影到第 $n$ 个平面波得到\eqnrefs{eq:kd_exact_momentum_chain}。在 $p_x=-\hbar\kappa_{\rm L}$ 时，$n=0,1$ 两态的自由能同为 $E_{\rm KD}$，而最近的外部态能量为 $9E_{\rm KD}$。将其余格点记作 $Q$ 空间，Feshbach/SW 展开给出
\[
H_{\rm eff}=PHP+PHQ\frac{1}{E-QH_0Q}QHP+\cdots.
\]
最低阶只保留 $PVP$，得到 $E_{\rm KD}\mathbb I_2+(U_p/4)\sigma_x$。最近两个外部态各给
\[
\delta E^{(2)}
=\frac{(U_p/4)^2}{E_{\rm KD}-9E_{\rm KD}}
=-\frac{U_p^2}{128E_{\rm KD}},
\]
在对称工作点对两个目标态完全相同，所以只产生共同相位。忽略 $O(U_p^2/\Delta_{\rm off})$ 差分修正后，去掉 $E_{\rm KD}\mathbb I_2$，二态传播子是
\[
U_2(\tau)=
\cos\!\left(\frac{U_p\tau}{4\hbar}\right)\mathbb I_2
-\ii\sin\!\left(\frac{U_p\tau}{4\hbar}\right)\sigma_x.
\]
对初态 $|-\hbar\kappa_{\rm L}\rangle$ 取另一分量模方，即得到\eqnrefs{eq:kd_bragg_rabi_probability}。
\end{exbox}

两束反向传播激光的强度条纹构成电子物质波光栅。短脉冲期间电子来不及分辨不同反冲能量，条纹只把空间依赖相位刻到波函数上；长脉冲则把自由色散变成谱选择器，只允许满足 Bragg 条件的近共振动量级持续交换。两种极限来自同一驻波动量链，差别只是实验解析的时间尺度；该动量链及其两个受控极限已经在上面的例题中逐步算出。

PINEM 与 Kapitza--Dirac 衍射都可能产生 Bessel 级次，但共同点只在于短时极限都化成“周期相位的 Fourier 系数”。它们的微观顶角与格点坐标不同：
\begin{center}
\begin{tabular}{lll}
\toprule
 & PINEM & Kapitza--Dirac\\
\midrule
主导顶点 & $A\!\cdot\!p$ & $A^2$\\
场结构 & 同步纵向场 & 驻波强度 $\cos^2(\kappa_{\rm L}x)$\\
格点坐标 & $E_0+\ell\hbar\omega$ & $p_x+n\hbar\kappa_{\rm KD}$\\
短时极限 & 能量-空间相位光栅 & 横向相位光栅\\
长时极限 & 反冲/谱选择 & Bragg 衍射\\
\bottomrule
\end{tabular}
\end{center}
